\medskip

\begin{tabularx}{\linewidth}{|X|}\hline
\textbf{Principaux domaines abordés :} Suites;
Fonctions, Fonction exponentielle.\\ \hline
\end{tabularx}

\bigskip

\textbf{Partie A}

\medskip

On considère la fonction $f$ définie pour tout réel $x$ par:

\[f(x) = 1+x - \text{e}^{0,5x - 2}.\]

On admet que la fonction $f$ est dérivable sur $\R$. On note $f'$ sa dérivée.

\medskip

\begin{enumerate}
\item 
	\begin{enumerate}
		\item Déterminer la limite de la fonction $f$ en $- \infty$.
		\item Démontrer que, pour tout réel $x$ non nul, $f(x) = 1+ 0,5x\left(2 - \dfrac{\text{e}^{0,5x}}{0,5x} \times \text{e}^{-2}\right)$.
		
En déduire la limite de la fonction $f$ en $+\infty$.
	\end{enumerate}
\item 
	\begin{enumerate}
		\item Déterminer $f'(x)$ pour tout réel $x$.
		\item Démontrer que l'ensemble des solutions de l'inéquation $f'(x) < 0$ est l'intervalle

$]4 + 2\ln (2)~;~+\infty[$.
	\end{enumerate}	
\item Déduire des questions précédentes le tableau de variation de la fonction $f$ sur $\R$. On fera figurer la valeur exacte de l'image de $4 + 2\ln (2)$ par $f$.
\item Montrer que l'équation $f(x) = 0$ admet une unique solution sur l'intervalle $[-1~;~0]$.
\end{enumerate}

\bigskip

\textbf{Partie B}

\medskip

On considère la suite $\left(u_n\right)$ définie par $u_0 = 0$ et, pour tout entier naturel $n$,
\: $u_{n+1} = f\left(u_n\right)$ où $f$ est la fonction définie à la partie A.

\medskip

\begin{enumerate}
\item 
	\begin{enumerate}
		\item Démontrer par récurrence que, pour tout entier naturel $n$, on a : 
		
		\[u_n \leqslant u_{n+1} \leqslant 4.\]
		
		\item En déduire que la suite $\left(u_n\right)$ converge. On notera $\ell$ la limite.
	\end{enumerate}	
\item 
	\begin{enumerate}
		\item On rappelle que $f$ vérifie la relation $\ell = f(\ell)$.
		
Démontrer que $\ell  = 4$.
		\item ~
		
\begin{minipage}{0.48\linewidth}On considère la fonction \texttt{valeur} écrite ci-contre dans le langage Python :
		\end{minipage}\hfill 
\begin{minipage}{0.48\linewidth}
\begin{tabular}{|l|}\hline
\texttt{def valeur (a) :}\\
\quad \texttt{u = 0}\\
\quad \texttt{n = 0}\\
\quad \texttt{while u $\leqslant$ a:}\\
\qquad \texttt{u=1 + u - exp(0.5*u - 2)}\\
\qquad \texttt{n = n+1}\\ 
\quad return \texttt{n}\\ \hline
\end{tabular}
\end{minipage}

L'instruction \texttt{valeur}(3.99) renvoie la valeur 12.

Interpréter ce résultat dans le contexte de l'exercice.
	\end{enumerate}
\end{enumerate}

\bigskip
